% RESEARCH_DRAFT_SHA256: 52471917a03f647c31af983bd453889921a56ddd96079aa6b8e844249abb4cac
% Layout and source-reference wrappers expanded for the web renderer.
\documentclass[11pt]{article}
\usepackage[margin=25mm]{geometry}
\usepackage{amsthm}
\usepackage{fontspec,unicode-math}
\setmainfont{Latin Modern Roman}
\setmathfont{Latin Modern Math}
\usepackage{amsmath,amssymb,mathtools,xurl,hyperref}
\hypersetup{colorlinks=true,urlcolor=blue,linkcolor=blue}
\setcounter{secnumdepth}{0}
\setlength{\parindent}{0pt}
\setlength{\parskip}{5pt}
\setlength{\emergencystretch}{3em}
\newcommand{\sref}[2]{\hyperlink{source:#1:#2}{[S#2]}}
\newcommand{\cataloguescope}[1]{\paragraph{Source record.}#1}
\newcommand{\floor}[1]{\left\lfloor#1\right\rfloor}
\newtheorem{proposition}{Proposition}
\newtheorem{lemma}[proposition]{Lemma}
\begin{document}
\paragraph{Research manuscript.} OpenAI \texttt{gpt-6-astra}, reasoning \texttt{ultra}; 11 October 2026.
\paragraph{Preserved source body.} The complete catalogue mathematical text and sources follow with layout and source links expanded, including source qualifications and its historical preparation notes. New research begins at the section headed Research attempt.

\section{A dual-fiber adequacy contract for the LaTeX ecosystem}\label{statement}

\subsection{Definitions and mathematical statement}

\paragraph{Shared notation.}Standard number systems are $\mathbb N,\mathbb Z,\mathbb Q,\mathbb R,\mathbb C$. Specialized notation and index ranges follow the cited source and the definitions below.


TeX is a typesetting system, and LaTeX supplies document structure and macros built on it. Mathematical source notation, document organization, visual layout, and accessibility are distinct authoring concerns. A replacement or supplement must be evaluated against stated workflows rather than an undefined overall notion of suitability.
\paragraph{Record scope and qualifications.}
The catalogue title is retained as metadata. Named certificates, contracts, bounds or frameworks in that title are not asserted as results of the original source extract.
The population of authors, requirements, interoperability constraints, and evaluation criteria are not specified. No particular replacement is endorsed.

Catalogue screening limitation: Workshop planning or software infrastructure agenda; the supplied passage does not identify a determinate mathematical target.

The source snapshot labels this record partially solved; this is catalogue metadata rather than independent verification of a mathematical result or the current completion of a project.
\paragraph{Statement recorded in the supplied source.}
This group should aim to understand to what extent LaTeX does and does not meet the needs of current mathematicians. After such an understanding is had, it should aim to create a list of resources that can supplement TeX, or replace it outright. This group may also venture into ideas on what tools can be developed to address points from the first objective and to create or plan to create said tools. \href{http://aimpl.org/cyberinfrastructure/10/}{[S2]}
\paragraph{Source status.}
The supplied catalogue labels this source record partially solved; that metadata was checked against the supplied snapshot on 2026-10-11. The label is preserved for provenance and is not an independent verification of present open status or mathematical correctness.
\subsection{Short English statement}
Assess which authoring needs LaTeX serves and identify useful supplements or alternatives.
\subsection{Sources}
\begin{itemize}
\item[S1] Ulam AI and the UnsolvedMath Contributors, UnsolvedMath: A Curated Collection of Open Mathematics Problems, catalogue entry. Catalogue curation and metadata: CC BY 4.0. The source extract is retained; definitions, English summaries and qualifications are editorial additions.
\url{https://huggingface.co/datasets/ulamai/UnsolvedMath}.
\item[S2] AIM Workshop Problem Lists, source linked by the catalogue.
\url{http://aimpl.org/cyberinfrastructure/10/}.
\end{itemize}
\paragraph{Preparation and review.}
The original source and UnsolvedMath contributors retain attribution for the source extract. Codex (GPT-6; exact serving revision not exposed) assisted with definitions, English summaries, scope notes and formatting. No new human mathematical authorship or independently verified solution is asserted. Catalogue identity checked against the complete supplied registry and current repository. Distinct upstream records, including repeated questions and broad or damaged extracts, are retained for complete coverage of the source CS category. A repeated record is not a new independent conjecture. Complete mathematical-equivalence and current-literature review has not been performed. Preparation and final review follow contribution prompts C and D, with source gaps disclosed. Structural and rendering checks do not establish a new theorem or open status.
\label{end:statement}

\section{Research attempt}
\subsection{Literature, source context and scope}
The complete catalogue body above retains the original agenda, source status, and historical preparation. The AIMPL problem page could not be retrieved. The supplied complete source transcription was read, and the accessible primary AIM workshop announcement was checked for context. The announcement is not a substitute for verification of the inaccessible problem page. The catalogue's earlier proposed frameworks and research summaries are context, not established solutions credited to this attempt. The following is an explicitly limited investigation of one route; it neither substitutes a formalized toy model for the entire source agenda nor claims empirical or deployed-software results. Elementary auxiliary facts are not new original-target progress.

\subsection{Testing whether visual mathematical output determines meaning}
The original agenda concerns what TeX supplies and what other tools might add. We isolate a semantic recovery route: infer the intended mathematical operation from a displayed expression alone. Let $f(x)=x^2$ on the positive reals and consider the displayed notation $f^{-1}(x)$. One author intends the inverse function; another intends the reciprocal of the function value. Their intended functions are respectively $\sqrt{x}$ and $1/x^2$.
\begin{proposition}
The same visible notation can encode different numerical functions. No decoder whose sole input is that notation can always recover the intended function in these two cases.
\end{proposition}
\begin{proof}
For $x>0$, the inverse of $f$ is $\sqrt{x}$ because squaring and taking the positive square root undo each other. The reciprocal value is $1/f(x)=1/x^2$. At $x=4$ these are 2 and $1/16$, so they differ. The visible input to the decoder is identical, hence its output is identical; it cannot equal both distinct intended functions. The same argument applies to any encoding map $E$ with $E(s_1)=E(s_2)$ for distinct semantic records: a left inverse $D$ would imply $s_1=D(E(s_1))=D(E(s_2))=s_2$, a contradiction.
\end{proof}
This is an ambiguity of notation without contextual conventions, not a claim that TeX itself promises to infer semantic intent or that inverse notation should normally mean reciprocal. A structured annotation identifying the operation would distinguish these two examples. It would not by itself meet every authoring, accessibility, archival, collaboration or rendering need.
\subsection{Verification and remaining gaps}
The two functions were evaluated at 4 with exact values. No needs assessment, resource list, replacement platform or comprehensive adequacy contract was delivered. The failed recovery route supplies a concrete reason to retain semantics where required, not an empirical verdict on the entire LaTeX ecosystem.

\section{References and source attribution}
Primary source: \url{http://aimpl.org/cyberinfrastructure/10/}. Attempted primary AIMPL page retrieval; failed with timeout or gateway error. The supplied full source transcription was read; the inaccessible live page was not independently inspected.

\url{https://aimath.org/pastworkshops/cyberinfrastructure.html}. Read the primary December 2023 workshop announcement; it describes the open-source mathematical infrastructure agenda.

All retrievals were checked on 11 October 2026. Original contributor and catalogue attribution remains in the preserved Sources section. The elementary arguments below are self-contained; their exposition is not a claim of first discovery.


\paragraph{Integration review.} Codex (GPT-6; exact serving revision not exposed) reviewed the written arguments and expanded layout and source-reference wrappers. Repository metadata and the production web renderer were checked using Python 3.13 and JavaScript with MathJax 3.2.2. This model-assisted review is not independent human proof verification; no PDF compilation or formal proof checking is asserted.

\section{Completion Estimate}
\noindent\textbf{COMPLETION: 0\%}
No substantiated new progress toward the entire original target is claimed. The worked results and failed route are delimited above; a known result, illustrative model, or auxiliary obstruction is not counted as original-target completion.

\section{Final status and attribution}
\textbf{Attempt status: unresolved.} This describes this research attempt, not a new ruling on the source catalogue's status or on external literature claims.
The mathematics and this self-audit were generated by OpenAI \texttt{gpt-6-astra}, reasoning \texttt{ultra}, on 11 October 2026. Codex (GPT-6; exact serving revision not exposed) orchestrated the batch and contributed formatting/repository review. No human mathematical authorship, independent human proof checking, or formal verification is claimed. The historical preparation paragraphs in the reproduced catalogue body concern that body only.

Contribution prompt A was applied to the complete definition and supplied source record, followed by the adversarial checks of prompt B. Local reproduction files are \texttt{batch013/checks.py}, \texttt{computations.json}, and the per-ID audit JSON. They were not public links at preparation time. Numerical checks are finite and do not prove unrestricted results. No LaTeX engine was found locally; compilation was not performed by this research agent.

\end{document}
